\section{ACF Earth Digital Twin}

The Atmospheric Complexity Framework implements a Digital Twin
approach for continuous monitoring, simulation and prediction of
the Earth atmosphere system.


\section{Digital Twin Architecture}

The ACF Digital Twin is composed of:

\begin{itemize}
\item Observation layer
\item Data assimilation layer
\item Simulation layer
\item Artificial intelligence layer
\item Visualization layer
\item Decision support layer
\end{itemize}



\section{Real-Time Earth System Synchronization}

The digital twin continuously integrates:

\begin{itemize}
\item Satellite observations
\item Radar observations
\item Ground stations
\item Ocean measurements
\item Numerical weather prediction outputs
\end{itemize}



\section{Data Flow Architecture}

The general workflow is:

\begin{equation}
Observation
\rightarrow
Assimilation
\rightarrow
Simulation
\rightarrow
Prediction
\rightarrow
Decision
\end{equation}



\section{High Performance Computing}

Large scale atmospheric simulation requires HPC resources.

ACF supports:

\begin{itemize}
\item Parallel computing
\item Distributed computing
\item GPU acceleration
\item Cloud computing
\item Exascale simulation concepts
\end{itemize}



\section{Model Acceleration}

The computational performance can be improved by:

\begin{itemize}
\item Domain decomposition
\item MPI parallelization
\item GPU kernels
\item AI surrogate models
\end{itemize}
